\documentclass[a4paper,12pt]{article}

\usepackage[T2A]{fontenc}
\usepackage[utf8]{inputenc}
\usepackage[russian]{babel}
\usepackage{amsmath,amssymb,amsthm,mathtools}
\usepackage{helvet}
\renewcommand{\familydefault}{\sfdefault}
\usepackage{icomma}
\usepackage{geometry}
\usepackage{microtype}
\usepackage{tikz}
\usepackage{tabularx,booktabs,longtable}
\usepackage{enumitem}
\usepackage{hyperref}
\usetikzlibrary{positioning}
\usepackage{parskip}
\usepackage{fancyhdr}
\usepackage{setspace}
\usepackage{xcolor}

\geometry{margin=2.1cm}
\onehalfspacing

\definecolor{accent}{HTML}{008080}
\definecolor{darkaccent}{HTML}{004D4D}
\definecolor{textgray}{HTML}{333333}
\definecolor{softgray}{HTML}{707070}
\definecolor{linegray}{HTML}{D8DEDE}

\hypersetup{
    colorlinks=true,
    linkcolor=darkaccent,
    urlcolor=darkaccent
}

\setlength{\parindent}{0pt}
\setlength{\parskip}{0.65em}

\setlist[itemize]{
    leftmargin=1.4em,
    itemsep=0.25em,
    topsep=0.3em
}

\setlist[enumerate]{
    leftmargin=1.8em,
    itemsep=0.65em,
    topsep=0.4em
}

\pagestyle{fancy}
\fancyhf{}
\fancyhead[L]{\small\textcolor{softgray}{Проценты}}
\fancyhead[R]{\small\textcolor{softgray}{Кирилл Зиновьев}}
\fancyfoot[C]{\small\textcolor{softgray}{\thepage}}
\renewcommand{\headrulewidth}{0.3pt}
\renewcommand{\headrule}{
    \hbox to\headwidth{\color{linegray}\leaders\hrule height \headrulewidth\hfill}
}

\newcommand{\pct}[1]{#1\,\%}

\begin{document}

% ============================================================
% ТИТУЛЬНЫЙ ЛИСТ
% ============================================================

\begin{titlepage}
\thispagestyle{empty}

\begin{center}

\vspace*{3.0cm}

{\Large\textcolor{darkaccent}{\textbf{ЧЕК-ЛИСТ}}}

\vspace{0.7cm}

{\huge\textbf{Проценты}}

\vspace{0.35cm}

{\Large
Три основных типа задач\\
и работа с пропорцией
}

\vspace{2.2cm}

{\large
\textcolor{textgray}{
Лёвин Артём Александрович\\
эксклюзивно для Кирилла Зиновьева
}}

\vspace{0.6cm}

{\normalsize 7 класс}

\vfill

{\large 26 августа 2026 г.}

\end{center}

\vspace{1.2cm}

\begin{tikzpicture}[remember picture,overlay]
    \draw[
        darkaccent,
        line width=1.2pt
    ]
    ([xshift=2.0cm,yshift=2.0cm]current page.south west)
    .. controls
    ([xshift=5.0cm,yshift=3.2cm]current page.south west)
    and
    ([xshift=-5.0cm,yshift=1.2cm]current page.south east)
    ..
    ([xshift=-2.0cm,yshift=2.3cm]current page.south east);
\end{tikzpicture}

\end{titlepage}

% ============================================================
% ОСНОВНОЙ ЧЕК-ЛИСТ
% ============================================================

\section*{\textcolor{darkaccent}{1. Главная идея}}

В задачах на проценты удобно переводить условие в пару соответствий:

\[
\boxed{\text{проценты} \longleftrightarrow \text{число}}
\]

Главная отправная точка:

\[
\boxed{\pct{100}\ \text{соответствует исходному числу}}
\]

Исходное число --- это число, \textbf{от которого} считают проценты.

Например:

\[
\pct{20}\text{ от }500.
\]

Здесь число \(500\) является исходным:

\[
500 \longleftrightarrow \pct{100}.
\]

Если исходное число неизвестно, напротив \(\pct{100}\) ставим \(x\).

\section*{\textcolor{darkaccent}{2. Универсальная таблица}}

Практически любую базовую задачу на проценты можно начать с такой схемы:

\begin{center}
\renewcommand{\arraystretch}{1.45}
\begin{tabular}{cc}
\toprule
\textbf{Проценты} & \textbf{Соответствующее число} \\
\midrule
\(\pct{100}\) & исходное число \\
\(\pct{p}\) & соответствующее значение \\
\bottomrule
\end{tabular}
\end{center}

Какая-либо из величин может быть неизвестна.

Например:

\[
\begin{array}{c|c}
\pct{100} & 600\\
\pct{30} & x
\end{array}
\]

или

\[
\begin{array}{c|c}
\pct{100} & x\\
\pct{30} & 600
\end{array}
\]

или

\[
\begin{array}{c|c}
\pct{100} & 50\\
x\,\% & 10
\end{array}
\]

Именно расположение \(x\) определяет тип задачи.

% ============================================================

\section*{\textcolor{darkaccent}{3. Как определить исходное число}}

Полезный ориентир:

\[
\boxed{\text{число, от которого считают проценты}
\longleftrightarrow \pct{100}}
\]

Особенно внимательно смотрите на слово \textbf{«от»}.

\subsection*{Пример}

Найдите \(\pct{30}\) от числа \(600\).

Фраза

\[
\pct{30}\text{ от числа }600
\]

показывает, что проценты считаются от \(600\).

Поэтому:

\[
\begin{array}{c|c}
\pct{100} & 600\\
\pct{30} & x
\end{array}
\]

\subsection*{Сравните}

Найдите число, если \(\pct{30}\) от него составляет \(600\).

Теперь число, от которого считаются проценты, неизвестно:

\[
\begin{array}{c|c}
\pct{100} & x\\
\pct{30} & 600
\end{array}
\]

Формулировки похожи, однако объект поиска различается.

% ============================================================

\section*{\textcolor{darkaccent}{4. Тип I. Найти процент от числа}}

Типовая формулировка:

\[
\boxed{\text{Найдите }\pct{p}\text{ от числа }a}
\]

Здесь исходное число известно:

\[
a \longleftrightarrow \pct{100}.
\]

Искомое значение соответствует \(\pct{p}\):

\[
x \longleftrightarrow \pct{p}.
\]

Получаем пропорцию:

\[
\frac{a}{100}=\frac{x}{p}.
\]

Перемножаем крест-накрест:

\[
100x=ap.
\]

Следовательно,

\[
\boxed{x=\frac{ap}{100}}.
\]

То же правило можно записать так:

\[
\boxed{\pct{p}\text{ от }a
=
a\cdot\frac{p}{100}}.
\]

\subsection*{Пример 1}

Найдите \(\pct{75}\) от \(600\).

\[
\begin{array}{c|c}
\pct{100} & 600\\
\pct{75} & x
\end{array}
\]

\[
\frac{600}{100}=\frac{x}{75}.
\]

\[
100x=600\cdot75.
\]

\[
x=\frac{600\cdot75}{100}=450.
\]

\[
\boxed{450}
\]

\subsection*{Пример 2}

Найдите \(\pct{20}\) от \(500\).

\[
x=\frac{500\cdot20}{100}=100.
\]

\[
\boxed{100}
\]

% ============================================================

\section*{\textcolor{darkaccent}{5. Тип II. Найти исходное число}}

Типовая формулировка:

\[
\boxed{
\text{Найдите число, если }\pct{p}\text{ этого числа составляет }b
}
\]

Теперь исходное число неизвестно:

\[
x \longleftrightarrow \pct{100}.
\]

Из условия известно:

\[
b \longleftrightarrow \pct{p}.
\]

Составляем пропорцию:

\[
\frac{x}{100}=\frac{b}{p}.
\]

Перемножаем крест-накрест:

\[
px=100b.
\]

Следовательно,

\[
\boxed{x=\frac{100b}{p}}.
\]

Также можно рассуждать через десятичную дробь:

\[
\boxed{x=b:\frac{p}{100}}.
\]

\subsection*{Пример 1}

Найдите число, если \(\pct{70}\) этого числа составляет \(371\).

\[
\begin{array}{c|c}
\pct{100} & x\\
\pct{70} & 371
\end{array}
\]

\[
\frac{x}{100}=\frac{371}{70}.
\]

\[
70x=371\cdot100.
\]

\[
x=\frac{371\cdot100}{70}=530.
\]

\[
\boxed{530}
\]

Проверка:

\[
530\cdot0{,}7=371.
\]

\subsection*{Пример 2}

Найдите число, если \(\pct{75}\) этого числа составляет \(600\).

\[
x=\frac{600\cdot100}{75}=800.
\]

\[
\boxed{800}
\]

Проверка:

\[
800\cdot0{,}75=600.
\]

% ============================================================

\section*{\textcolor{darkaccent}{6. Тип III. Найти, сколько процентов составляет одно число от другого}}

Типовая формулировка:

\[
\boxed{
\text{Какой процент составляет число }a\text{ от числа }b?
}
\]

Число после слова \textbf{«от»} является исходным:

\[
b \longleftrightarrow \pct{100}.
\]

Числу \(a\) соответствует неизвестное количество процентов:

\[
a \longleftrightarrow x\,\%.
\]

Пропорция:

\[
\frac{b}{100}=\frac{a}{x}.
\]

Перемножаем крест-накрест:

\[
bx=100a.
\]

Следовательно,

\[
\boxed{x=\frac{100a}{b}}.
\]

Или короче:

\[
\boxed{x=\frac{a}{b}\cdot100\%}.
\]

\subsection*{Пример 1}

Какой процент составляет \(10\) от \(50\)?

\[
\begin{array}{c|c}
\pct{100} & 50\\
x\,\% & 10
\end{array}
\]

\[
\frac{50}{100}=\frac{10}{x}.
\]

\[
50x=1000.
\]

\[
x=20.
\]

Следовательно,

\[
\boxed{\pct{20}}.
\]

\subsection*{Пример 2}

Какой процент составляет \(10\) от \(80\)?

\[
x=\frac{10}{80}\cdot100=12{,}5.
\]

\[
\boxed{\pct{12,5}}.
\]

Десятичные проценты являются обычным результатом задачи.

% ============================================================

\section*{\textcolor{darkaccent}{7. Три формулировки рядом}}

\renewcommand{\arraystretch}{1.45}

\begin{table}[h]
\centering
\caption{Как меняется положение неизвестной величины}
\begin{tabularx}{\linewidth}{>{\raggedright\arraybackslash}X
                             >{\raggedright\arraybackslash}X}
\toprule
\textbf{Формулировка} & \textbf{Что обозначаем через \(x\)} \\
\midrule

Найдите \(\pct{25}\) от числа \(50\).
&
Число, соответствующее \(\pct{25}\):
\[
\pct{100}\leftrightarrow50,\qquad
\pct{25}\leftrightarrow x.
\]
\\

Найдите число, \(\pct{25}\) которого равны \(50\).
&
Исходное число:
\[
\pct{100}\leftrightarrow x,\qquad
\pct{25}\leftrightarrow50.
\]
\\

Какой процент составляет \(10\) от \(80\)?
&
Количество процентов:
\[
\pct{100}\leftrightarrow80,\qquad
x\,\%\leftrightarrow10.
\]
\\

\bottomrule
\end{tabularx}
\end{table}

Ответы для сравнения:

\[
\pct{25}\text{ от }50=12{,}5;
\]

\[
\pct{25}\text{ некоторого числа}=50
\quad\Longrightarrow\quad
x=200;
\]

\[
10\text{ от }80
\quad\Longrightarrow\quad
\pct{12,5}.
\]

% ============================================================

\section*{\textcolor{darkaccent}{8. Алгоритм решения через пропорцию}}

Для базовых задач на проценты используйте один и тот же порядок действий.

\begin{enumerate}
    \item Найдите в условии число, \textbf{от которого} считают проценты.
    \item Поставьте это число напротив \(\pct{100}\).
    \item Если исходное число неизвестно, поставьте напротив \(\pct{100}\) букву \(x\).
    \item Во второй строке запишите указанный процент и соответствующее ему число.
    \item Неизвестную величину обозначьте через \(x\).
    \item Составьте пропорцию.
    \item Перемножьте крест-накрест.
    \item Решите получившееся уравнение.
    \item Проверьте, что единицы ответа соответствуют вопросу:
    число или проценты.
\end{enumerate}

% ============================================================

\section*{\textcolor{darkaccent}{9. Быстрая схема выбора}}

\begin{center}
\begin{tikzpicture}[
    every node/.style={font=\small},
    block/.style={
        draw=darkaccent,
        rounded corners=2pt,
        inner xsep=9pt,
        inner ysep=7pt,
        text width=8.0cm,
        align=center
    },
    arrow/.style={
        ->,
        >=stealth,
        line width=0.8pt,
        darkaccent
    }
]

\node[block] (start)
{\textbf{Что требуется найти?}};

\node[block, below=8mm of start] (type1)
{%
Значение \(\pct{p}\) от известного числа \(a\)\\[3pt]
\(\displaystyle x=\frac{ap}{100}\)
};

\node[block, below=8mm of type1] (type2)
{%
Исходное число, если известно значение его \(\pct{p}\)\\[3pt]
\(\displaystyle x=\frac{100b}{p}\)
};

\node[block, below=8mm of type2] (type3)
{%
Сколько процентов составляет \(a\) от \(b\)\\[3pt]
\(\displaystyle x=\frac{100a}{b}\)
};

\draw[arrow] (start.south) -- (type1.north);
\draw[arrow] (type1.south) -- (type2.north);
\draw[arrow] (type2.south) -- (type3.north);

\end{tikzpicture}
\end{center}

% ============================================================

\section*{\textcolor{darkaccent}{10. Типичные ошибки}}

\subsection*{Ошибка 1. Считать любое число в условии процентами}

Запись

\[
10
\]

означает обычное число.

Запись

\[
\pct{10}
\]

означает десять процентов.

Символ \(\%\) принципиально важен.

\subsection*{Ошибка 2. Неверно выбрать значение, соответствующее \(\pct{100}\)}

Всегда спрашивайте:

\[
\boxed{\text{От какого числа считают процент?}}
\]

Именно это число соответствует \(\pct{100}\).

\subsection*{Ошибка 3. Ставить \(x\) напротив привычной строки}

\(x\) ставится напротив той величины, которую требуется найти.

Если требуется найти проценты, то:

\[
x\,\%
\]

находится в столбце процентов.

Если требуется найти число, \(x\) находится в столбце числовых значений.

\subsection*{Ошибка 4. Неверно перемножать пропорцию}

Если

\[
\frac{a}{b}=\frac{c}{d},
\]

то

\[
\boxed{ad=bc}.
\]

Перемножаются величины, расположенные крест-накрест.

\subsection*{Ошибка 5. Пугаться десятичного ответа}

Например,

\[
\pct{2}\text{ от }10
=
10\cdot\frac{2}{100}
=
0{,}2.
\]

Ответ \(0{,}2\) полностью корректен.

Также допустимы дробные значения процентов:

\[
\pct{12,5},\qquad
\pct{2,4},\qquad
\pct{37,5}.
\]

% ============================================================

\section*{\textcolor{darkaccent}{11. Самопроверка перед ответом}}

Перед тем как завершить задачу, проверьте четыре пункта:

\begin{itemize}
    \item Я правильно определил число, соответствующее \(\pct{100}\)?
    \item Я поставил \(x\) напротив той величины, которую требуется найти?
    \item Я правильно перемножил пропорцию крест-накрест?
    \item Мой ответ является именно тем, о чём спрашивали: числом или процентами?
\end{itemize}

\bigskip

\begin{center}
\textcolor{darkaccent}{
\large\textbf{
Главное: сначала определите, что является \(\pct{100}\),\\
и только затем составляйте пропорцию.
}}
\end{center}

% ============================================================
% ТРЕНИРОВОЧНЫЙ БЛОК
% ============================================================

\newpage

\section*{\textcolor{darkaccent}{Тренировочный блок}}

\textbf{Путь Б.}

Решайте каждую задачу через таблицу соответствий или пропорцию.

Для каждой задачи сначала определите:

\[
\boxed{\text{что соответствует }\pct{100}}
\]

и только после этого выполняйте вычисления.

% ------------------------------------------------------------

\subsection*{\textcolor{darkaccent}{Среда, 26 августа 2026 г.}}

\begin{enumerate}
    \item Найдите \(\pct{20}\) от числа \(350\).

    \item Найдите число, если \(\pct{25}\) этого числа составляет \(60\).

    \item Какой процент составляет число \(18\) от числа \(60\)?

    \item Найдите число, если \(\pct{40}\) этого числа составляет \(140\).

    \item Найдите \(\pct{15}\) от числа \(260\).
\end{enumerate}

\vspace{0.7cm}

\textcolor{softgray}{
Подсказка дня: число после слова «от» чаще всего помогает определить,
что соответствует \(\pct{100}\).
}

% ------------------------------------------------------------

\subsection*{\textcolor{darkaccent}{Четверг, 27 августа 2026 г.}}

\begin{enumerate}
    \item Какой процент составляет число \(24\) от числа \(80\)?

    \item Найдите \(\pct{35}\) от числа \(240\).

    \item Найдите число, \(\pct{12}\) которого составляют \(42\).

    \item Найдите \(\pct{2}\) от числа \(45\).

    \item Какой процент составляет число \(45\) от числа \(120\)?
\end{enumerate}

\vspace{0.7cm}

\textcolor{softgray}{
Подсказка дня: если требуется найти проценты,
неизвестная величина \(x\) должна находиться в столбце процентов.
}

% ------------------------------------------------------------

\subsection*{\textcolor{darkaccent}{Пятница, 28 августа 2026 г.}}

\begin{enumerate}
    \item Найдите число, если \(\pct{18}\) этого числа составляет \(81\).

    \item Какой процент составляет число \(14\) от числа \(40\)?

    \item Найдите \(\pct{7,5}\) от числа \(320\).

    \item Найдите число, \(\pct{65}\) которого составляют \(247\).

    \item Какой процент составляет число \(27\) от числа \(72\)?
\end{enumerate}

\vspace{0.7cm}

\textcolor{softgray}{
Подсказка дня: дробный или десятичный ответ может быть полностью корректным.
Проверяйте смысл результата, а не только его внешний вид.
}

% ============================================================
% ОТВЕТЫ
% ============================================================

\newpage

\section*{\textcolor{darkaccent}{Ответы к тренировочному блоку}}

\renewcommand{\arraystretch}{1.45}

\begin{table}[h]
\centering
\caption{Среда, 26 августа 2026 г.}
\begin{tabularx}{0.82\linewidth}{c X}
\toprule
\textbf{№} & \textbf{Ответ} \\
\midrule
1 & \(70\) \\
2 & \(240\) \\
3 & \(\pct{30}\) \\
4 & \(350\) \\
5 & \(39\) \\
\bottomrule
\end{tabularx}
\end{table}

\vspace{0.8cm}

\begin{table}[h]
\centering
\caption{Четверг, 27 августа 2026 г.}
\begin{tabularx}{0.82\linewidth}{c X}
\toprule
\textbf{№} & \textbf{Ответ} \\
\midrule
1 & \(\pct{30}\) \\
2 & \(84\) \\
3 & \(350\) \\
4 & \(0{,}9\) \\
5 & \(\pct{37,5}\) \\
\bottomrule
\end{tabularx}
\end{table}

\vspace{0.8cm}

\begin{table}[h]
\centering
\caption{Пятница, 28 августа 2026 г.}
\begin{tabularx}{0.82\linewidth}{c X}
\toprule
\textbf{№} & \textbf{Ответ} \\
\midrule
1 & \(450\) \\
2 & \(\pct{35}\) \\
3 & \(24\) \\
4 & \(380\) \\
5 & \(\pct{37,5}\) \\
\bottomrule
\end{tabularx}
\end{table}

\vfill

\begin{center}
\textcolor{darkaccent}{
\textbf{
Сначала определите \(\pct{100}\).\\
Затем расположите данные.\\
После этого пропорция становится почти механической.
}}
\end{center}

\end{document}